\section{Estimation and use at an unseen location}
\label{sec:use}

\paragraph{Prediction at an unseen location.}
Algorithm~1 of \cite{ng2025} applies with environments $u=(y,t,e)$. A VAE with a shared encoder and
decoder is trained on all environments. The prior over $\hat Z_c$ is a Gaussian ANM whose mechanism
parameters are conditioned on $c=(y,t)$ (for instance through embeddings of $y$ and $t$), and the prior over
$\hat Z_s$ has mean and variance conditioned on $r=(t,e)$. The adjacency matrix of the content block is learned
with the sparsity and moral-graph penalties of \cite{ng2025}, and sinks are identified and frozen
iteratively. The style coordinates may be designated in advance as isolated nodes or discovered as estimated
nodes with no edges whose distribution depends on $(t,e)$ but not on the species $y$; an isolated content
latent with no parents whose distribution depends on $y$ is a content latent. The numbers of latents $n_c$
and $n_s$ are hyperparameters; \cite{ng2025} suggest choosing the number of latents by the ELBO.
(\Cref{app:code} documents, against the released reference code, which of these pieces are already
implemented, which are only partly implemented, and which still need to be built for M1.)

\paragraph{Prediction at a new location $e^\ast$.}
Encode $x$ to $(\hat z_c,\hat z_s)$. Since $x\leftrightarrow(z_c,z_s)$ is an injection, the 
posterior of $C=(Y,T)$ is
$$
p(c\mid x^\ast,e^\ast) = p(c\mid z_c,z_s,e^\ast) = p(y,t\mid z_c,z_s,e^\ast) =
p(y\mid t,z_c,z_s,e^\ast)p(t\mid z_c,z_s,e^\ast).
$$
Two cases exist: Time of day observed or unobserved at test time.
\begin{itemize}
\item \emph{Time of day observed at test time.} Since $Z_s$ is independent of $(Y,Z_c)$
	given $(x^\ast,t^\ast,e^\ast)$,
	\begin{equation}
	p(y\mid x^\ast,t^\ast,e^\ast) =
	p(y\mid z_c,z_s,t^\ast,e^\ast) =
	p(y\mid z_c,t,e^\ast)\ \propto\ p^{e^\ast}(y\mid t^\ast)\,p(z_c\mid y,t^\ast),
	\label{eq:predict}
	\end{equation}
	so the style latents drop out exactly and only the content latents are used.
	Hence, the species $y$ can be inferred from maximizing the likelihood-prior product over $y$ for 
	the given $t^\ast$.
	The factor $p(z_c\mid y,t^\ast)$ is invariant by \labelcref{item:P2}.
	The prior $p^{e^\ast}(y\mid t^\ast)$ can be obtained from the location-specific prior 
	$p^{e^\ast}(C)$;
	it can be taken equal to the pooled training prior, or estimated from unlabeled images of 
	$e^\ast$ (for
	instance by EM), if multiple images are available, or assumed to be uniform and ignored. No 
	style parameters of $e^\ast$ are needed.

\item \emph{Time of day not observed.} Then
	$p(c\mid z_c,z_s,e^\ast)\propto p^{e^\ast}(c)\,p(z_c\mid c)\,p^{(t,e^\ast)}(z_s)$, and the last 
	factor depends on $t$, so the style carries information about the time of day. Dropping it gives
	$p(c\mid z_c,e^\ast)\propto p^{e^\ast}(c)\,p(z_c\mid c)$, which is the correct posterior given 
	$z_c$
	alone but not the full posterior; using the style would require the style distributions of the 
	new  location for each time of day, which are not available from labelled data.

	Coarse time of day could be approximated by the image's brightness:
	This is reasonable, with one caveat. Brightness is mostly determined by style (exposure, IR 
	mode) and by $t$, so it works as a cheap, fixed read-out of part of $z_s$. It fails if a new 
	camera shifts the day/night threshold, for example a flash-lit night image or a shaded day 
	image.

	We can measure this before relying on it. Fit brightness $\rightarrow$ day/night on some on the 
	training locations and test on the rest. Then there are two cases:
	\begin{enumerate}
	\item If the error is near zero, treat $t$ as observed and use
		$p(y \mid z_c, t) \ \propto\ \bar{p}(y|\tilde{t})p(z_c \mid y,\tilde{t})$, where $\tilde{t}$ 
		is the estimated $t$ from brightness.
	\item If not, use it softly:
		$p(y \mid z_c, b) \ \propto\ \sum_t \bar{p}(y,t)p(z_c \mid y,t)p(b \mid t)$.
		This assumes brightness $b$ depends on the image only through $t$ and style, not on the 
		species content, which holds roughly unless animal coat color moves overall brightness.

		The soft version still needs $p(b \mid t)$ at L100, which you take from the pooled training 
		data. That is the same unknown style distribution as before, but one-dimensional, so a rough 
		version is far more benign than the full style vector.
\end{enumerate}
\end{itemize}
